\section{The splitting criterion}%
\label{sec:asym_splitting}

Periodic data alone does not in general imply the existence of sitewise
fusion tensors (Remark~\ref{rem:asym_zipper}). The following criterion gives
a sufficient algebraic condition for the compressions of
Theorem~\ref{thm:asym_multiblock} to split into sitewise intertwiners.

\begin{theorem}[Splitting criterion]
    \label{thm:asym_splitting}
    \lean{MPSTensor.exists_multiBlockCompression_remainder_eq_zero_of_star,
        MPSTensor.exists_multiBlockCompression_remainder_eq_zero_of_isSemisimpleModule,
        MPSTensor.isSemisimpleModule_wordModule_of_conjTranspose_mem_adjoin}
    \leanok
    \uses{thm:asym_multiblock, def:asym_word_module}
    In the setting of Theorem~\ref{thm:asym_multiblock}, suppose that the
    unital algebra generated by the matrices $B^i$ is closed under conjugate
    transposition, that is, $(B^i)^\dagger$ lies in the unital algebra generated
    by the $B^j$ for every $i$. Then the word-algebra module of $B$ is
    semisimple, and $B$ admits a block triangular form as in
    Theorem~\ref{thm:asym_multiblock} whose remainder vanishes: the block
    triangular form is block diagonal.
\end{theorem}

\begin{proof}
    \leanok
    \uses{thm:asym_multiblock, thm:asym_flag_theorem}
    Since the algebra is closed under conjugate transposition, the
    orthogonal complement of an invariant subspace is invariant. Thus
    every invariant subspace has an invariant complement and the module
    is semisimple. Every step of the flag of
    Theorem~\ref{thm:asym_flag_theorem} then has an invariant complement,
    the module is the direct sum of these complements, and a basis
    adapted to the direct sum makes every $B^i$ block diagonal with the
    blocks prescribed by the flag.
\end{proof}

\begin{corollary}[Sitewise fusion tensors]
    \label{cor:asym_fusion_tensors_star}
    \lean{MPSTensor.exists_fusionTensors_of_star}
    \leanok
    \uses{thm:asym_splitting}
    Under the hypotheses of Theorem~\ref{thm:asym_splitting} there are
    maps $V_s$ and $W_s$ with $W_sV_t=\delta_{st}\id$, $W_sB^wV_s=C_s^w$
    for every word $w$, and the sitewise relations $B^iV_s=V_sC_s^i$ and
    $W_sB^i=C_s^iW_s$ for every physical index $i$ and block occurrence $s$.
\end{corollary}

\begin{center}
    % Formula: B^iV_s=V_sC_s^i and W_sB^i=C_s^iW_s for every letter i.
    % Source: Corollary~\ref{cor:asym_fusion_tensors_star}.
    % Ink-to-index map: the ring is the compression map $V_s$ or $W_s$;
    % the square node is $B$ or its target $C_s$.
    % Contraction set: the single virtual bond joining the ring to the
    % square node in each pair; the physical leg $i$ and the two outer
    % virtual legs remain open.
    % Boundary signature: (virtual-west=1 | virtual-east=1 |
    % physical-up=1 | physical-down=0) for each pictured equality.
    \tenkzkernel
    \begin{tenkz}[rows={wire}, cols=2, west=open, east=open]
        \tn[label pos=s, ports={90:physical:$i$}]{B} & \tn[skin=ring]{V_s}
    \end{tenkz}
    \quad $=$ \quad
    \begin{tenkz}[rows={wire}, cols=2, west=open, east=open]
        \tn[skin=ring]{V_s} & \tn[label pos=s, ports={90:physical:$i$}]{C_s}
    \end{tenkz}
    \\[2ex]
    \begin{tenkz}[rows={wire}, cols=2, west=open, east=open]
        \tn[skin=ring]{W_s} & \tn[label pos=s, ports={90:physical:$i$}]{B}
    \end{tenkz}
    \quad $=$ \quad
    \begin{tenkz}[rows={wire}, cols=2, west=open, east=open]
        \tn[label pos=s, ports={90:physical:$i$}]{C_s} & \tn[skin=ring]{W_s}
    \end{tenkz}
\end{center}

\begin{proof}
    \leanok
    \uses{thm:asym_splitting}
    With vanishing remainder, $B^i=\sum_sV_sC_s^iW_s$, and the sitewise
    relations follow from the biorthogonality of the pairs.
\end{proof}

\begin{corollary}[Obstruction to closure under conjugate transposition]
    \label{cor:asym_anomaly_not_star}
    \lean{MPSTensor.not_forall_conjTranspose_mem_adjoin_of_forall_right_intertwiner_eq_zero,
        CZXCompression.czxSquare_not_starClosed}
    \leanok
    \uses{thm:asym_splitting}
    If, in the setting of Theorem~\ref{thm:asym_multiblock}, some block
    occurrence $s$ admits no nonzero sitewise right intertwiner $X$ with
    $B^iX=XC_s^i$ for all $i$, then the unital algebra generated by the
    $B^i$ is not closed under conjugate transposition. In particular, the
    stacked product of the anomalous CZX symmetry with itself, for which
    both sitewise intertwiner spaces vanish, generates a unital algebra
    that is not closed under conjugate transposition.
\end{corollary}

\begin{proof}
    \leanok
    \uses{thm:asym_splitting}
    By Theorem~\ref{thm:asym_splitting}, closure under conjugate
    transposition would give a compression pair whose right map is a
    nonzero sitewise intertwiner.
\end{proof}

\begin{remark}
    \label{rem:asym_anomaly_unitary}
    The physical operator of the CZX symmetry, a product of controlled-Z
    gates and single-site spin flips, is unitary at every length. The
    failure of closure under conjugate transposition in
    Corollary~\ref{cor:asym_anomaly_not_star} therefore concerns the
    virtual algebra of the stacked tensor only, not the physical operator;
    unitarity of the operator is not part of the formal statement.
    This algebraic conclusion does not compute the anomaly class.
    Determining that class requires a separate computation of the
    $3$-cocycle defined by the associator of the fusion tensors.
\end{remark}

\section{The single-block reduction of Moln\'ar--Ge--Schuch--Cirac}%
\label{sec:asym_single_block}

\begin{corollary}[Single-block case: the normal-target reduction]
    \label{cor:asym_single_block}
    \lean{MPSTensor.exists_isReduction_and_nilpotent_of_isNormal}
    \leanok
    \uses{def:same_mpv2_pos, def:normal, def:mps_rectangular_reduction}
    Let $A$ be a normal tensor (NT) of bond dimension $D_A\geq1$, let $B$
    be an arbitrary tensor of bond dimension $D_B$, and suppose that
    $\ket{V^{(N)}(B)}=\ket{V^{(N)}(A)}$ for every $N\geq1$.
    Then there are maps $V:\C^{D_A}\to\C^{D_B}$ and
    $W:\C^{D_B}\to\C^{D_A}$ with $WV=\id$ such that, for every $N\geq0$
    and every word of length $N$,
    \begin{align}
        WB^{i_1}\cdots B^{i_N}V & =A^{i_1}\cdots A^{i_N},
        \label{eq:asym_single_block_compression}
    \end{align}
    and the remainder $R^i=B^i-VA^iW$ generates a nilpotent non-unital
    algebra. Explicitly, for every word of length $r=D_B-D_A+1$,
    \begin{align}
        R^{i_1}R^{i_2}\cdots R^{i_r} & =0.
        \label{eq:asym_single_block_nilpotency}
    \end{align}
    In particular $D_B\geq D_A$.
\end{corollary}

\begin{center}
    % Formula: WB^{i_1}\cdots B^{i_N}V=A^{i_1}\cdots A^{i_N}
    % (\ref{eq:asym_single_block_compression}).
    % Source: Corollary~\ref{cor:asym_single_block}.
    % Ink-to-index map: the two rings are the single compression maps
    % $W$ (left) and $V$ (right), the special case of $W_s$, $V_s$ at
    % the unique slot; the interior nodes are copies of $B$.
    % Contraction set: every horizontal virtual bond between consecutive
    % copies of $B$ and between each $B$ and its adjacent ring; the
    % physical legs and the two outer virtual legs (the
    % $D_A$-dimensional indices of the matrix $A^{i_1}\cdots A^{i_N}$)
    % remain open.
    % Boundary signature: (virtual-west=1 | virtual-east=1 |
    % physical-up=2 | physical-down=0) for the finite schematic shown,
    % with the two explicit lettered nodes carrying the drawn physical
    % ports.
    \tenkzkernel
    \begin{tenkz}[rows={wire}, cols=5, west=open, east=open]
        \tn[skin=ring]{W} &
        \tn[label pos=s, ports={90:physical:$i_1$}]{B} &
        \tn[skin=dots]{} &
        \tn[label pos=s, ports={90:physical:$i_N$}]{B} &
        \tn[skin=ring]{V}
    \end{tenkz}
    \quad $=$ \quad
    \begin{tenkz}[rows={wire}, cols=3, west=open, east=open]
        \tn[label pos=s, ports={90:physical:$i_1$}]{A} &
        \tn[skin=dots]{} &
        \tn[label pos=s, ports={90:physical:$i_N$}]{A}
    \end{tenkz}
\end{center}

\begin{proof}
    \leanok
    \uses{thm:asym_multiblock}
    Apply Theorem~\ref{thm:asym_multiblock} to the single-block target
    $C=A$. Clauses (iv) and (v) give the compression pair and
    (\ref{eq:asym_single_block_compression}), clause (vii) gives
    $D_B=D_A+z\geq D_A$, and clause (vi) gives
    (\ref{eq:asym_single_block_nilpotency}) with $r=|S|+z=1+(D_B-D_A)$.
\end{proof}

The proportional case, in which $\ket{V^{(N)}(B)}=\lambda^N\ket{V^{(N)}(A)}$
for a scalar $\lambda\in\C^\times$, reduces to
Corollary~\ref{cor:asym_single_block} by replacing $B$ with
$\lambda^{-1}B$, and then $WB^{i_1}\cdots B^{i_N}V=\lambda^NA^{i_1}\cdots A^{i_N}$
with remainder $R^i=B^i-\lambda\,VA^iW$; the compression pair of that case
is recorded in Theorem~\ref{thm:mps_proportional_reduction_exists}.

Corollary~\ref{cor:asym_single_block} together with the proportional case
corresponds to Propositions~20 and~21 of~\cite{Molnar2018SemiInjective} (with
$V,W$ swapped), with three source corrections and an explicit nilpotency bound.
First, the hypothesis $\lambda\ne0$ is required because the printed
Proposition~20 fails at $\lambda=0$: the zero tensor $B=0$ satisfies the
hypothesis with no compression onto a nonzero normal target. Second, the
source proof silently normalizes $\lambda=1$; the reduction to the equality case
makes the scalar explicit. Both corrections are recorded in
\cite{gap:mgsc18_reduction_proportionality_scalar}
(\texttt{docs/paper-gaps/mgsc18_reduction_proportionality_scalar.tex}). Third, the remainder
algebra is generated without unit (a unital algebra contains a nonzero identity
and is not nilpotent). The triangular gauge of Theorem~\ref{thm:asym_multiblock}
yields the explicit nilpotency length bound $r=D_B-D_A+1$; see
\cite{gap:mgsc18_nilpotency_length_one_terminology}
(\texttt{docs/paper-gaps/mgsc18_nilpotency_length_one_terminology.tex}). The uniqueness statement
(Theorem~22 of~\cite{Molnar2018SemiInjective}) is not addressed here and fails
verbatim in the multi-block setting, as the matching isomorphisms $X_s$ of
Theorem~\ref{thm:asym_multiblock} are not unique in general.
